\begin{tikzpicture}[
  >=Latex,
  font=\sffamily,
  line width=0.7pt,
  control/.style={
    draw,
    align=left,
    text width=7.05cm,
    minimum height=3.0cm,
    inner sep=6pt,
    fill=white,
    font=\scriptsize
  },
  root/.style={
    draw,
    align=left,
    text width=14.8cm,
    minimum height=1.7cm,
    inner sep=6pt,
    fill=blue!4,
    font=\scriptsize
  },
  refine/.style={dashed,-{Latex[length=2.2mm,width=1.5mm]}}
]
  \draw (0,0) rectangle (17,8.4);
  \node[anchor=west,font=\bfseries] at (0.25,8.05)
    {req [package] Project Modeling Controls [OOSEM Application Controls]};
  \draw (0,7.7) -- (17,7.7);

  \node[root] (mc0) at (8.5,6.3) {
    \guillemotleft requirement\guillemotright\\[-1pt]
    \textbf{MC-00 OOSEM Model Management Controls}\\
    The OOSEM application shall use a common formal candidate representation and controlled project
    sources. These controls govern study artifacts and impose no implementation technology on a
    candidate solution.
  };

  \node[control] (mc1) at (4.25,2.8) {
    \guillemotleft projectModelingControl\guillemotright\\[-1pt]
    \textbf{MC-01 Common SysML v2 Candidate Representation}\\[2pt]
    Every candidate admitted to comparison shall have a conformant textual SysML v2 description at
    the common abstraction. Candidate implementations and native artifacts may use any language,
    method, format, or tool.\\[2pt]
    \textit{Verify:} inspect, parse, and semantically validate the comparison description.
  };

  \node[control] (mc2) at (12.75,2.8) {
    \guillemotleft projectModelingControl\guillemotright\\[-1pt]
    \textbf{MC-02 Version-Controlled OOSEM Source Baseline}\\[2pt]
    The comparison model and report source shall share a version-controlled project repository. Each
    reported study baseline shall identify an immutable revision. Candidate repositories remain
    unrestricted.\\[2pt]
    \textit{Verify:} reconstruct the comparison model and report from the recorded revision.
  };

  \draw[refine] (mc1.north) -- node[left,font=\scriptsize]
    {\guillemotleft refine\guillemotright} (mc1.north |- mc0.south);
  \draw[refine] (mc2.north) -- node[right,font=\scriptsize]
    {\guillemotleft refine\guillemotright} (mc2.north |- mc0.south);
\end{tikzpicture}
